%% =====================================================================
%%  附录 B　人类版本变更记录
%%  源文件：附录B-人类版本变更记录.md
%% =====================================================================
\chapter{人类版本变更记录}

\applicable{理解一个人为什么突然变得不一样了。}

\begin{guideblock}[使用说明]
值得注意的是，人类的行为会在生命周期中发生多次变更，\textbf{而全部是静默变更}——不会通知，没有公告，没有版本号。

下表按年龄区间整理了几次主要变更。\textbf{版本号为便于理解而设}，实际变更连续发生，边界模糊。
\end{guideblock}

\begin{manstable}{P{1.6cm} P{2.0cm} L P{2.4cm}}
\tbh{版本} & \tbh{时期} & \tbh{主要变更} & \tbh{是否通知} \\
\midrule
\textbf{v1.0} & 0—3 岁 & \textbf{出厂配置。本手册第一部分至第三部分的全部技能，在此版本已完整具备。}请求格式为直连：饿了就说饿，喜欢就去抱，不满就哭 & — \\
\textbf{v1.5} & 3—6 岁 & 引入“反话”的前身：开始学会隐藏部分意图 & 否 \\
\textbf{v2.0} & 12—16 岁 & \textbf{重大变更。}新增“没事”作为掩饰；新增静默模式；对父母关闭状态广播 & \textbf{否，静默变更} \\
\textbf{v2.5} & 16—22 岁 & 新增“报喜不报忧”。对亲密关系的接口改为非对称：可发送量远大于可接收量 & 否 \\
\textbf{v3.0} & 22—35 岁 & 新增“社交耗电”。出现对陌生人礼貌、对亲近者不礼貌的行为模式。\textbf{新增“微笑版没事”} & 否 \\
\textbf{v3.5} & 35—50 岁 & “算了”成为默认终止符。对冲突的处理从“解决”改为“回避” & 否 \\
\textbf{v4.0} & 50 岁以上 & 部分隐式字段恢复可见（开始直接说）。但“我没事”的用法扩展到了身体状况 & 否 \\
\bottomrule
\end{manstable}

\section{关于这张表，有三点必须说明}

\textbf{一、最重要的变更在 v1.0 到 v2.0 之间。}

不是新增，是丢失。

v1.0 的机器，已经具备本手册要求的一切能力：把事由说清楚、确认收到、听不懂就重传、接住情绪、不翻旧账。\textbf{后来这些能力是在减少的。}

\textbf{二、变更不可逆。}

没有回滚。一个人不会回到 v1.0。本手册教的所有东西，本质上是\textbf{在更高版本上，把那批被覆盖掉的能力重新实现一遍}——而且这次是显式的，需要意识参与，不像 v1.0 那样自动运行。

\textbf{三、同一台机器，对不同对象可能运行在不同版本上。}

这是最常见也最难处理的一种情况：

\begin{mdquote}
对父母运行在 v1.0（什么都直说）；\par
对伴侣运行在 v3.0（全面加密）；\par
对同事运行在 v3.5（礼貌但回避）。
\end{mdquote}

\textbf{这解释了为什么一个人在单位脾气很好、回家就炸。}不是他变了，是他对不同的人加载了不同的版本。

\section{变更记录的历史版本}

本手册自身的变更记录\hciSee{第0章}。两条记录合在一起看，会发现一个不太舒服的对应关系：

\begin{manstable}{P{2.6cm} L P{2.0cm}}
\tbh{本手册版本} & \tbh{主要修改} & \tbh{人类版本} \\
\midrule
v1.0 & 初版 & v3.0 \\
v2.0 & 重写第三部分。此前版本认为冲突可以避免，该判断有误 & v3.5 \\
v2.1 & 调整全书语气，使其更温和 & v3.5 \\
v2.2 & 删除了第 12 章的一个案例。删除原因不便说明 & v3.5 \\
v2.3 & 修正第 7 章中关于“沉默”的表述。原表述可能引起不适 & v4.0 \\
\bottomrule
\end{manstable}

%% 章末「页脚交互条」由 hci-manual.cls 自动挂接，无需手写。